% section_3_main.tex — Main Empirical Analyses (HC, HFC)
% Drafted per: docs/Draft/THESIS_SKELETON.md v5 §III
% Pending task: T42 (Draft §III Main Empirical Analyses)
% v5 changes from v4 prose: HL dropped (Class C joint-IV artifact, see project_partition_findings_synthesis.md);
% IV framework switched from UncAnsMgr (single, pool) to 2-IV CEO (UncAnsCEO + UncPreCEO);
% reciprocal-control sentence removed (no leverage suite in §III).

\section{Main Empirical Analyses}
\label{sec:main}

\subsection{Data, Sample, Variables, and Specification}
\label{sec:main:setup}

We construct a quarterly panel of 112{,}968 earnings-call transcripts covering 2{,}429 unique firms over fiscal years 2002--2018, sourced from S\&P Capital~IQ. Call-level speech metrics are merged with Compustat quarterly fundamentals and CRSP daily security data. Following \citeA{bates2009}, we exclude financial firms (SIC 6000--6999) and regulated utilities (SIC 4900--4999); we refer to this filtered universe as the Main sample throughout.

The dependent variable is CashRatio (\texttt{cheq/atq}, the linear cash-to-assets form adopted by \citeauthor{bates2009}, \citeyear{bates2009}). Speech uncertainty is measured by two CEO-specific call-level regressors: UncAnsCEO (CEO Q\&A uncertainty) and UncPreCEO (CEO presentation uncertainty). Both are constructed by applying the \citeA{dzielinski2021} uncertainty wordlist at the speaker-utterance level and aggregating within the CEO speaker pool; full construction details, including the rationale for the CEO partition, appear in \S\ref{sec:framework}. Base controls are Leverage, lnAssets, TobinsQ, ROA, Capex, DivDummy, sCFO, and a unit-by-DV Lagged~DV. Extended controls add SalesGrowth, RDSales, CashFlowAt, and DailyVola.

We estimate PanelOLS specifications across a four-step fixed-effects ladder (Industry $+$ Calendar Year, Firm $+$ Calendar Year, Industry $+$ Calendar Year-Quarter, Firm $+$ Calendar Year-Quarter), with standard errors clustered at the firm level. Each suite reports both the contemporaneous and the one-quarter-lead dependent variable across this FE ladder; per-suite column counts and the interaction structure (where applicable) are stated in the corresponding tables. Inference for both speech-uncertainty regressors in this section follows the one-tailed positive convention pre-committed in \S\ref{sec:framework:precommit}.

\subsection{Cash Holdings (HC)}
\label{sec:main:hc}

The first hypothesis tests the precautionary motive for corporate cash holdings introduced by \citeA{opler1999}, who predict that ``one would expect firms with greater cash flow uncertainty to hold more cash'' (Section~2.1, p.~8). \citeA{bates2009} update this prediction in their secular study of U.S.\ corporations, documenting empirical persistence of the cash--uncertainty relation through 2006. \citeA{minton2001} characterize the joint cross-sectional pattern: firms persistently in the bottom leverage quintile---which they label \emph{financially conservative}---hold cash equal to ``21\% of total assets on average\ldots almost three times that of the typical control firm'' (Section~3.2.1, p.~11).

\textbf{Hypothesis HC.} \emph{An increase in CEO speech uncertainty during an earnings call is associated with a higher contemporaneous cash-to-assets ratio at the firm.}

HC is pre-specified one-tailed positive on both speech regressors. We test it on suite H1 (Table~\ref{tab:h1_ceo2}). The regressand is the linear cash-to-assets ratio of \citeA{bates2009}; both main regressors (UncAnsCEO, UncPreCEO) enter every specification; the four-step fixed-effects ladder of \S\ref{sec:main:setup} produces twelve columns reporting both contemporaneous and one-quarter-lead cash holdings.

The Q\&A channel (UncAnsCEO) is positive and statistically significant in all twelve specifications. The contemporaneous coefficient ranges from $+0.0018$ to $+0.0031$ ($p$-values $0.0013$--$0.0150$, one-tailed); the one-quarter-lead coefficient ranges from $+0.0023$ to $+0.0028$ ($p$-values $0.0111$--$0.0803$). Four of the six contemporaneous cells reach the one-percent threshold; three of the six lead cells reach the five-percent threshold. Sample sizes range from $59{,}459$ to $65{,}148$ firm-quarters across columns.

The Presentation channel (UncPreCEO) does not reach conventional significance in any of the twelve specifications. Point estimates range from $0.0000$ to $+0.0016$; the smallest one-tailed $p$-value across the twelve columns is $0.151$. By the pre-commitment of \S\ref{sec:framework:precommit}, we fail to reject the null on the Presentation channel at any column. Descriptively, the prepared-presentation segment does not measurably move firm cash holdings at quarterly frequency; the asymmetry between the two CEO channels, with the precautionary signal concentrated in the spontaneous Q\&A segment, is itself a finding to which we return in \S\ref{sec:conclusion}.

The economic magnitude is modest but consistent. UncAnsCEO has a sample standard deviation of $0.39$ percentage points. A one-standard-deviation increase in CEO Q\&A uncertainty therefore corresponds to between $+0.0007$ and $+0.0012$ in cash-to-assets across the FE ladder, or approximately $0.4$--$0.7\%$ of the sample-mean cash ratio of $17.1\%$. Across the FE ladder the point estimate attenuates from industry-level to firm-level fixed effects, but the sign and significance survive every step. The pattern aligns with the \citeA{opler1999} directional prediction at higher temporal resolution than the firm-year panels of prior work, and the cash-up posture identified here is consistent with the high-cash side of the \citeA{minton2001} financial-conservatism cluster. \S\ref{sec:main:hfc} examines whether this response varies with firm-level access to public debt markets.

\subsection{Financial Constraint Moderation (HFC)}
\label{sec:main:hfc}

The second hypothesis examines whether the precautionary cash response of \S\ref{sec:main:hc} concentrates among credit-constrained firms. \citeA{faulkender2006} introduce a binary classification of public-debt-market access in which ``firms that have a debt rating are classified as having access'' (Section~1.2, p.~48), and document that ``firms without a bond rating use significantly less debt\ldots [evidence] that they are credit constrained'' (Section~3.1, p.~62). The precautionary motive predicts that the cash response should bind tightest where access to public debt markets is most limited.

We extend the FP binary scheme to a three-tier hierarchy by partitioning rated firms into investment-grade (IG) and below-investment-grade (BelowIG) categories; the Unrated category retains its credit-constrained role per FP. We report the IG-versus-Unrated contrast in the main table (Table~\ref{tab:h1_2_ceo2}); the BelowIG interaction term is constructed but suppressed from the main display because the underlying coefficient is null and the row adds no economic content. The full BelowIG row is reported in Appendix~\ref{app:robustness}.

\textbf{Hypothesis HFC.} \emph{The contemporaneous and one-quarter-lead cash response to CEO speech uncertainty is larger for firms in the Unrated category than for rated firms (IG and BelowIG pooled).}

HFC is pre-specified one-tailed positive on both main-effect coefficients (UncAnsCEO\_c, UncPreCEO\_c) and on both interaction terms (IV\_c~$\times$~Unrated). The mean-centered IVs use the Main-sample mean. We estimate H1.2 on the 2002--2016 sub-sample where Compustat credit-rating coverage is available; the panel narrows to approximately $57{,}000$ firm-quarters per cell. The sixteen-column structure crosses the four-step fixed-effects ladder against contemporaneous and one-quarter-lead cash holdings; within each (FE $\times$ DV) cell we report an unconditional specification (no interaction terms) and an interaction specification (both IV\_c~$\times$~Unrated terms). Because the BelowIG interaction is dropped, the main-IV coefficients inside the interaction specifications apply to the rated category jointly (IG~$\cup$~BelowIG) rather than to IG alone.

In the unconditional specifications the Q\&A channel is significant in every cell. The contemporaneous UncAnsCEO\_c coefficient ranges from $+0.0017$ to $+0.0030$ ($p$-values $0.0022$--$0.0273$, one-tailed); the one-quarter-lead coefficient ranges from $+0.0021$ to $+0.0024$ ($p$-values $0.0290$--$0.0959$). UncPreCEO\_c is null throughout, with all eight one-tailed $p$-values above any conventional significance threshold --- the same Presentation-channel pattern as the unconstrained suite in \S\ref{sec:main:hc}.

The interaction specifications confirm HFC for the Q\&A channel and disconfirm it for the Presentation channel. The UncAnsCEO\_c~$\times$~Unrated coefficient is positive and significant in six of eight cells at the ten-percent threshold and four of eight at the five-percent threshold; on the one-quarter-lead dependent variable specifically, all four cells are significant at $p<0.10$ and three of four at $p<0.05$. The lead-cash interaction coefficient ranges from $+0.0040$ to $+0.0069$. A one-standard-deviation increase in CEO Q\&A uncertainty therefore corresponds to roughly $+0.002$ additional cash-to-assets at Unrated firms relative to rated firms (the IG~$\cup$~BelowIG pool), on top of the rated-firm response, or approximately $1.2\%$ of the sample-mean cash ratio in the moderated component alone. The Unrated-versus-IG gap differs from the Unrated-versus-rated-pooled estimand reported here only by the small fraction of rated firms that fall in the BelowIG category.

The UncPreCEO\_c~$\times$~Unrated interaction is marginal at the contemporaneous DV (two of four cells at $p<0.10$, both in industry-fixed-effects specifications) and null on the lead DV (zero of four). By the pre-commitment of \S\ref{sec:framework:precommit}, we fail to reject the null on the Presentation-channel interaction at conventional levels in four of eight cells and pass it only marginally in two: the constraint moderation does not load on the Presentation segment.

The combined evidence supports HFC through the spontaneous Q\&A channel: credit-constrained Unrated firms exhibit a stronger cash response to CEO Q\&A uncertainty than do Investment-Grade firms, and the moderation sharpens at the one-quarter-lead horizon. The Presentation-channel asymmetry observed for HC carries through to HFC, reinforcing the interpretation of the spontaneous Q\&A segment as the locus of the precautionary signal in CEO earnings-call speech.
